\section{Exact proof certificates and edge cases}
\label{app:proof-details}

This appendix expands the finite calculations used in the main proof.  The
purpose is auditability: none of these computations depends on a numerical
search or on a finite surrogate for an all-prime statement.

\subsection{Source-only periodic collapse}

There is also a direct source-only proof that \(0^{\mathbb Z}\) is the sole
periodic source point.  Suppose \(x\in\Xsf\) has period \(n\geq1\) and some
\(x_a=1\).  Choose a rational prime \(p\nmid n\).  Periodicity gives
\[
  a+n\mathbb Z\subseteq\supp(x).
\]
Because \(n\) is invertible modulo \(p^2\), the progression
\(a+n\mathbb Z\) meets every residue class modulo \(p^2\).  Therefore
\(\supp(x)\bmod p^2=\mathbb Z/p^2\mathbb Z\), contradicting admissibility.
Thus \(x=0^{\mathbb Z}\).

This argument alone is not enough for \cref{thm:factor-rigidity}: source
aperiodicity need not descend to a factor.  The CRT proof supplies the
stronger proximality property that does descend.

\subsection{CRT certificate for a finite window}

For fixed \(L\), enumerate the index set
\[
  I_L=[-L,L]\times\{1,2\}.
\]
Choose an injective map \((j,i)\mapsto p_{j,i}\) from \(I_L\) into the
rational primes.  The exact congruence packet is
\[
  n\equiv a_{j,i}-j\pmod{p_{j,i}^2}
  \qquad((j,i)\in I_L),
\]
where \(a_{j,i}\in M_{p_{j,i}}(x^{(i)})\).  Since distinct primes have
coprime squares, the Chinese remainder theorem gives one class
\[
  n\pmod{\prod_{(j,i)\in I_L}p_{j,i}^2}.
\]
No compatibility condition between the residues is needed.  Selecting the
least nonnegative representative produces a forward time \(n_L\geq0\).
For each \((j,i)\), the congruence says that \(n_L+j\) lies in a residue
missing from \(\supp(x^{(i)})\), and hence the indicated coordinate is zero.

The smallest-window case \(L=0\) uses two distinct primes \(p_x,p_y\) and
two congruences
\[
  n\equiv a_x\pmod{p_x^2},
  \qquad
  n\equiv a_y\pmod{p_y^2}.
\]
It synchronizes the central symbols.  Larger windows use the same argument
with fresh primes, not an induction that reuses potentially incompatible
moduli.

\subsection{Product-metric normalization}

For the metric in \cref{eq:product-metric}, the total mass is
\[
  \frac13\sum_{k\in\mathbb Z}2^{-|k|}
  =\frac13\left(1+2\sum_{k=1}^{\infty}2^{-k}\right)=1.
\]
If \(x_k=y_k\) for \(|k|\leq L\), then \(|x_k-y_k|\leq1\) outside the
window and
\[
  d_X(x,y)
  \leq\frac13\left(2\sum_{k=L+1}^{\infty}2^{-k}\right)
  =\frac{2^{1-L}}{3}.
\]
The right-hand side tends to zero.  This supplies an explicit convergence
rate, although the theorem needs only convergence.

\subsection{Why a periodic orbit has a positive gap}

Let \(y\) have least period \(r>1\).  If
\(S^k y=S^{k+1}y\) for some \(k\), applying \(S^{-k}\) gives \(y=Sy\),
contradicting least period \(r>1\).  Thus every one of the finitely many
numbers
\[
  d_Y(S^k y,S^{k+1}y),\qquad 0\leq k<r,
\]
is positive.  Their minimum \(\delta_r\) is positive.  For
\(n=qr+k\), periodicity gives
\[
  d_Y(S^n y,S^n(Sy))
  =d_Y(S^k y,S^{k+1}y)\geq\delta_r.
\]
This proves nonproximality of the pair \((y,Sy)\) without any expansivity or
symbolic structure.

For a second fixed point \(y\neq y_0\), the constant distance
\(d_Y(y,y_0)>0\) gives the analogous contradiction.  Both cases are needed:
the adjacent-orbit argument degenerates when \(r=1\).

\subsection{Fixed-point exponential and primitive accounting}

Substituting \(\#\Fix(S^m)=1\) into the definition yields
\[
  \log\zeta_{\AM,Y}(z)=\sum_{m=1}^{\infty}\frac{z^m}{m}.
\]
Formal differentiation gives
\[
  \frac{d}{dz}\log\zeta_{\AM,Y}(z)
  =\sum_{m=1}^{\infty}z^{m-1}=\frac{1}{1-z},
\]
and the constant term is zero, so
\(\log\zeta_{\AM,Y}(z)=-\log(1-z)\).  Exponentiating yields
\((1-z)^{-1}\).  Analytically the same calculation holds for \(|z|<1\).

The primitive Euler form has one factor because the only primitive orbit is
\(\mathcal O_0\):
\[
  \prod_{\mathcal O\in\operatorname{Prim}(Y,S)}
  (1-z^{|\mathcal O|})^{-1}
  =(1-z)^{-1}.
\]
Expanding the logarithm of this one factor produces all powers \(z^r\).
The expansion changes traversal number, not primitive identity.

\subsection{Finite-P0 witness, including the empty case}

For nonempty finite \(P_0\), \cref{eq:finite-witness} has support
\(1+Q\mathbb Z\).  Modulo any \(p^2\mid Q\), this support is exactly the
single residue \(1\), so every other residue is missing.  If \(d>0\) is a
period, the implication
\[
  x_1=1\Longrightarrow x_{1+d}=1
\]
forces \(1+d\equiv1\pmod Q\), hence \(Q\mid d\).  Since \(Q\) itself is a
period, the least period is exactly \(Q\).

When \(P_0=\varnothing\), the empty product is \(Q=1\) and admissibility
places no restriction on the full binary shift.  The all-zero and all-one
points are two distinct fixed points.  This is the correct edge-case
certificate; describing it as a nontrivial period-one orbit would be false.

